\section{Is My Variable Mapping to the ACC Domain Correct?}

\begin{contextbox}{The Tracking Control System Model}
	The system model from the paper is:
	\begin{equation}
		x(k+1) = x(k) + \tau\bigl[f(x(k)) + g(x(k))\, u(k)\bigr]
		\tag{System Model}
	\end{equation}
	
	With the designed control function, the deployed tracking equation becomes:
	\begin{equation}
		x(k+1) = x(k) - \beta\,\text{sgn}(x(k) - r(k)) + \Delta r(k)
		\tag{Tracking Eq.}
	\end{equation}
	
	where $\Delta r(k) = r(k+1) - r(k)$ and the error dynamics satisfy:
	\begin{equation}
		e(k+1) = e(k) - \beta\,\text{sgn}(e(k))
		\tag{Error Dynamics}
	\end{equation}
\end{contextbox}

\begin{questionbox}{Question: Does this variable mapping hold?}
	\renewcommand{\arraystretch}{1.5}
	\begin{center}
		\begin{tabular}{|c|l|l|}
			\hline
			\textbf{Variable} & \textbf{General Meaning}     & \textbf{ACC Interpretation}               \\
			\hline
			$r(k)$            & Reference trajectory         & Expected speed at time step $k$           \\
			\hline
			$\Delta r(k)$     & Change in reference          & Expected change in speed                  \\
			\hline
			$x(k)$            & System state                 & Reported (sensor) speed at time step $k$  \\
			\hline
			$e(k)$            & Tracking error $r(k) - x(k)$ & Discrepancy: expected vs. reported speed  \\
			\hline
			$u(k)$            & Control input                & Corrective action applied to reduce error \\
			\hline
			$g(x(k))$         & Input coupling function      & How the correction maps to state change   \\
			\hline
		\end{tabular}
	\end{center}
\end{questionbox}

\begin{mathbox}{Analysis: What Needs Clarification}
	The mapping is \textbf{mostly correct}, but there are nuances to think through:
	
	\begin{enumerate}[label=\textbf{\arabic*.}]
		\item \textbf{$r(k)$ --- ``Expected speed''}: This is the reference trajectory that the vehicle \emph{should} follow. In the ACC context, this comes from the upper controller's target speed profile (based on the driver's set speed and the safe following distance). This is well-defined.
		      
		\item \textbf{$x(k)$ --- ``Reported speed''}: Here is where the attack matters. Under normal operation, $x(k)$ is the true vehicle speed. Under a spoofing attack, $x(k)$ becomes the \emph{corrupted} sensor reading. The question is: does the tracking controller operate on the reported (possibly spoofed) value, or does it have access to a separate estimate of the true speed?
		      
		      \begin{itemize}
			      \item If the controller only sees the spoofed $x(k)$, then $e(k) = r(k) - x(k)$ is \emph{itself corrupted}, and the controller will issue wrong corrections.
			      \item If the controller has an independent estimate of true speed (e.g., from wheel encoders, IMU integration, or the IDS's detection), then $x(k)$ represents this estimate, and the tracking controller corrects against the attack.
		      \end{itemize}
		      
		      \textbf{This distinction is critical to define clearly in the paper.}
		      
		\item \textbf{$u(k)$ --- ``Correction function''}: In the general system model, $u(k)$ is the full control function:
		      \[
			      u(k) = g(x(k))^{-1}\bigl(-f(x(k)) + \Delta r(k) - \beta\,\text{sgn}(x(k) - r(k))\bigr)
		      \]
		      In the ACC analog, this represents the acceleration/deceleration command that the tracking controller issues to correct the speed error. This is analogous to (but distinct from) the PID's $u(t)$.
		      
		\item \textbf{$g(x(k))$ --- ``How correction affects reported speed''}: In the general model, $g(x(k))$ is the input coupling --- it determines how effectively the control input influences the state dynamics. In the ACC context, this captures the vehicle's response characteristics (e.g., how throttle/brake commands translate to speed changes, which depends on current speed, road grade, etc.).
	\end{enumerate}
\end{mathbox}
